\chapter{Einführungen}
\quellseite{11}

Im Auftrag des Ehrw. Präses Becker wurde P. Rud. M. Norden am 1. Sonnt. n. Epiph. in seiner Gemeinde zu Fremont, Nebr., eingeführt von Aug. Leuthäuser.

Am 22. Sonnt. n. Trin. wurde Lehrer E. Munzel in der Zions-Gemeinde zu Roseland (Chicago), Ill., eingeführt von Alfred E. Reinke.

Am Sonntag nach Neujahr wurde Lehrer E. M. Brandt als Lehrer der ersten Schulklasse der Bethel-Gemeinde zu Detroit, Mich., eingeführt von J. M. Gugel.

\quellseite{12}

Am Sonntag nach Neujahr wurde Lehrer Albert Winter als Lehrer der dritten Klasse der Dreieinigkeits-Gemeinde in Sheboygan, Wis., eingeführt von Fr. Mohlbrecht.

Am Sonntag nach Neujahr wurde Lehrer Georg Windisch als erster Lehrer an der Schule der Bethelem-Gemeinde (103. Straße) zu Chicago, Ill., eingeführt von H. Engelbrecht jun.

Am Sonntag nach Neujahr wurde Lehrer H. Klinger, berufen an die erste Klasse der Emmaus-Schule zu Indianapolis, Ind., eingeführt von Theo. Schurbel.

Am Sonntag nach Neujahr wurde Lehrer L. Fuhrmann als Lehrer an der Mittelklasse der St. Pauli’s-Negermissionsschule zu New Orleans, La., eingeführt von F. J. Lankenau.
